%=====================================================================
\chapter{Survey Research}\label{ch:survey}
%=====================================================================
\locator{3}{Part~3 --- Research Designs for Computing}

\begin{outcomes}
By the end of this chapter you will be able to:
\begin{tight}
  \item \textbf{Separate} population, sampling frame and achieved sample, and
    say what each limits.
  \item \textbf{Choose} a sampling method and state honestly what inference it
    supports.
  \item \textbf{Diagnose} nonresponse bias, and explain why response rate alone
    does not measure it.
  \item \textbf{Write} items free of the six common wording faults.
  \item \textbf{Decide} whether to adopt a validated instrument or build one.
  \item \textbf{Report} a survey to a recognised standard.
\end{tight}
\end{outcomes}

\begin{prereq}
\chref{questions} for constructs and operationalisation --- a questionnaire is
an operationalisation, and most survey failures are measurement failures.
\chref{measurement} follows this chapter and completes it.
\end{prereq}

\begin{keyterms}
population $\bullet$ sampling frame $\bullet$ coverage error $\bullet$
probability sampling $\bullet$ stratified sampling $\bullet$ cluster sampling
$\bullet$ convenience sampling $\bullet$ snowball sampling $\bullet$ response
rate $\bullet$ nonresponse bias $\bullet$ Likert item and Likert scale
$\bullet$ double-barrelled question $\bullet$ social desirability bias
$\bullet$ cognitive interview $\bullet$ pilot
\end{keyterms}

\section{What a Survey Can and Cannot Do}

A survey measures what a population reports about itself at one point in time.
That is genuinely valuable --- attitudes, practices, perceptions and self-reported
behaviour are often exactly what you need to know, and no other method gets
them at scale.

It is also a narrow instrument, and three limits are worth stating at the
outset.

\begin{center}
\small
\begin{tabular}{@{}L{4.4cm}L{9.9cm}@{}}
\toprule
\textbf{A survey cannot} & \textbf{Because} \\
\midrule
Establish causation
  & Nothing is manipulated and nothing is randomised. Cross-sectional survey
    data supports association (\chref{quasiexp}) \\
Measure behaviour
  & It measures \emph{reported} behaviour. People misremember, round, and
    present themselves favourably. Where a trace exists, mine it
    (\chref{quasiexp}) and treat the survey as complementary \\
Rescue a vague construct
  & If ``code quality'' is undefined, asking 300 people to rate it produces
    300 different constructs averaged into one number
    (\chref{measurement}) \\
\bottomrule
\end{tabular}
\end{center}

\section{Population, Frame, Sample}

\begin{definitionbox}[Three different things, routinely conflated]
The \term{population} is who you want to say something about --- ``professional
software developers in Pakistan''.

\medskip
The \term{sampling frame} is the list you can actually draw from --- members of
three professional groups on a messaging platform. The gap between population
and frame is \term{coverage error}, and it is usually the largest error in a
computing survey, larger than sampling error by a wide margin.

\medskip
The \term{sample} is who you invited; the \term{achieved sample} is who
responded. The gap between them is nonresponse.
\end{definitionbox}

\begin{alertbox}[The confidence interval does not know about coverage error]
A survey of 400 developers recruited from one online community can report a
margin of error of about five percentage points. That number describes sampling
error \emph{within the frame} only. It says nothing about the developers who
are not in that community, who may be the majority and may differ
systematically.

\medskip
Reporting the interval without discussing the frame implies a precision the
study does not have. State the frame, state who it excludes, and say in which
direction that would bias each main result.
\end{alertbox}

\section{Sampling}

\begin{center}
\small
\begin{longtable}{@{}L{2.9cm}L{5.0cm}L{6.2cm}@{}}
\toprule
\textbf{Method} & \textbf{How} & \textbf{Inference it supports} \\
\midrule
\endfirsthead
\toprule
\textbf{Method} & \textbf{How} & \textbf{Inference it supports} \\
\midrule
\endhead
\bottomrule
\endfoot
Simple random
  & Every frame member equally likely
  & To the frame, with quantified uncertainty. Rarely achievable in our field \\
Stratified
  & Divide the frame into strata, sample within each
  & As above, with better precision, and lets you report by stratum ---
    valuable when a subgroup is small but important \\
Cluster
  & Sample groups, then units within them
  & As above, but with a design effect: correlated units inside clusters mean
    the effective sample size is smaller than the count \\
Systematic
  & Every $k$th member
  & As simple random, unless the frame has a periodicity matching $k$ \\
\midrule
Convenience
  & Whoever responds
  & \emph{To nobody beyond the respondents.} Descriptive only \\
Purposive
  & Deliberately chosen for a characteristic
  & Analytic, not statistical (\chref{casestudy}) \\
Snowball
  & Respondents recruit others
  & Reaches hidden populations; biased towards the well-connected \\
Quota
  & Fill preset category counts non-randomly
  & Looks representative on the quota variables, is not on anything else \\
\end{longtable}
\end{center}

\begin{conceptbox}[Non-probability sampling is not disqualifying; pretending otherwise is]
Almost every software engineering survey uses convenience or snowball sampling,
because no sampling frame for ``software developers'' exists. That is a real
constraint, not a personal failing.

\medskip
What separates a good such survey from a bad one is the write-up. A good one
says: this is a convenience sample recruited from these three channels; here is
the respondent profile in detail; here is who is likely under-represented; the
results describe these respondents and generalise only as far as the reader
judges them typical. A bad one reports a margin of error and calls the sample
representative.

\medskip
\reg{PhD Regs 2024}{\S14.3(B)} asks about external validity --- ``can the
results be generalized to other settings and to other populations?'' The honest
answer for a convenience sample is ``to a limited degree, and here is why'',
which is an acceptable answer. A dishonest ``yes'' is not.
\end{conceptbox}

\section{Nonresponse}

\begin{alertbox}[Response rate is a weak proxy for the thing you actually care about]
The standard reflex is to worry about response rate. What matters is
\term{nonresponse bias} --- whether responders differ from non-responders
\emph{on the variables you are measuring}.

\medskip
These come apart in both directions. A 20\% response rate with responders and
non-responders alike on the measured variables produces an unbiased estimate. A
70\% response rate where the 30\% who declined are precisely the developers
who found the topic uncomfortable produces a badly biased one. Response rate is
worth reporting, but it does not settle the question.
\end{alertbox}

\begin{center}
\small
\begin{tabular}{@{}L{4.2cm}L{10.1cm}@{}}
\toprule
\textbf{Check} & \textbf{How} \\
\midrule
Compare with frame characteristics
  & If you know the frame's distribution on any variable, compare it with your
    respondents' \\
Wave analysis
  & Compare early responders with late ones. Late responders resemble
    non-responders more closely; a trend across waves suggests bias \\
Follow up a sample of non-responders
  & Even 20 short interviews with people who declined is far more informative
    than any rate \\
Report the funnel
  & Invited, opened, started, completed, usable. Each drop-off is a place bias
    can enter \\
\bottomrule
\end{tabular}
\end{center}

\section{Writing Items That Measure Something}

\begin{center}
\small
\begin{longtable}{@{}L{2.9cm}L{4.6cm}L{6.6cm}@{}}
\toprule
\textbf{Fault} & \textbf{Example} & \textbf{Why it breaks} \\
\midrule
\endfirsthead
\toprule
\textbf{Fault} & \textbf{Example} & \textbf{Why it breaks} \\
\midrule
\endhead
\bottomrule
\endfoot
Double-barrelled
  & ``Our tests are fast and reliable''
  & A respondent who thinks they are fast but flaky cannot answer \\
Leading
  & ``How much has our tool improved your workflow?''
  & Presupposes improvement; no option expresses ``it did not'' \\
Vague quantifier
  & ``Do you review code regularly?''
  & Regularly means daily to one person and monthly to another. Ask for a
    frequency \\
Assumes knowledge
  & ``Do you use mutation testing?''
  & Respondents who do not know the term will guess rather than admit it \\
Socially loaded
  & ``Do you write unit tests for your code?''
  & There is a right answer and everyone knows it. Ask what proportion of last
    week's commits had tests \\
Recall-heavy
  & ``How many bugs did you fix last year?''
  & Nobody knows. Shorten the window or use a trace \\
\end{longtable}
\end{center}

\subsection{Likert items, and one distinction that matters}

A \term{Likert item} is a single statement with an ordered agreement response.
A \term{Likert scale} is several items summed to measure one construct. The
distinction matters because it determines what analysis is legitimate: a single
item is ordinal, and a summed multi-item scale is closer to interval.
\chref{measurement} and \chref{inference} take this up; the design consequence
is that \textbf{if you intend to treat a construct as continuous, you must
write several items for it}, not one.

\begin{center}
\small
\begin{tabular}{@{}L{3.4cm}L{10.9cm}@{}}
\toprule
Points          & 5 or 7. More points do not add information once respondents
                  stop discriminating \\
Midpoint        & Include one unless you have a specific reason to force a
                  choice; forcing produces arbitrary answers, not decisive ones \\
Labels          & Label every point, not just the ends \\
Direction       & Reverse-word a few items to detect straight-lining --- but
                  avoid negations, which are misread \\
Don't know      & Offer it separately from the midpoint. They mean different
                  things and conflating them corrupts the mean \\
\bottomrule
\end{tabular}
\end{center}

\begin{conceptbox}[Adopt a validated instrument if one exists]
If you are measuring technology acceptance, usability, workload or job
satisfaction, validated instruments exist --- UTAUT2~\cite{venkatesh2012consumer},
SUS, NASA-TLX and others. Use them, unmodified, and cite them.

\medskip
The gain is substantial: the construct validity work is done, you can compare
your results with the literature, and reviewers do not have to take your word
for the measurement. Modifying items --- even lightly, even to fix awkward
wording --- discards the validation, so if you must adapt, say so explicitly and
re-establish reliability and validity on your own data (\chref{measurement},
\chref{sem}).
\end{conceptbox}

\begin{templatebox}[C8 --- Survey instrument checklist]
Work through this before the pilot, not after.

\medskip
\begin{tabular}{@{}L{0.8cm}L{13.0cm}@{}}
$\square$ & Every item traces to a construct, and every construct to a research
            question \\
$\square$ & Multi-item scales for any construct to be treated as continuous \\
$\square$ & No double-barrelled, leading or vague-quantifier items \\
$\square$ & Sensitive items behaviourally anchored rather than asked directly \\
$\square$ & Response options exhaustive and mutually exclusive \\
$\square$ & ``Don't know'' distinct from the midpoint \\
$\square$ & Demographics at the end, and only those actually analysed \\
$\square$ & Estimated completion time, measured on the pilot not guessed \\
$\square$ & Consent screen: purpose, storage, withdrawal, contact
            (\chref{ethics}) \\
$\square$ & Anonymous or identifiable --- decided, stated, and consistent with
            the data handling plan \\
$\square$ & Cognitive interviews with 3--5 people from the target population \\
$\square$ & Full pilot, with the analysis script run on pilot data \\
$\square$ & Instrument to be published as an appendix
            (\chref{reproducibility}) \\
\end{tabular}
\end{templatebox}

\subsection{Cognitive interviewing}

Sit with three to five people from the target population and have them complete
the questionnaire while thinking aloud. Ask what each question meant to them
and how they arrived at the answer.

This is the single highest-value hour in survey design. It reliably reveals
items that are read in a sense you never intended --- and unlike a pilot, which
tells you the distribution of answers, it tells you \emph{why} people answered
as they did.

\section{Reporting}

\begin{center}
\small
\begin{tabular}{@{}L{4.0cm}L{10.3cm}@{}}
\toprule
Population and frame  & Who you wanted, what list you had, who is missing \\
Recruitment           & Every channel, with dates \\
The funnel            & Invited, started, completed, usable, with exclusions
                        and reasons \\
Respondent profile    & Enough that a reader can judge typicality \\
Instrument            & In full, as an appendix or a deposit \\
Reliability           & For each multi-item scale (\chref{measurement}) \\
Analysis              & Including how ordinal items were treated and why \\
Nonresponse           & Which checks you ran and what they showed \\
\bottomrule
\end{tabular}
\end{center}

For internet surveys the CHERRIES checklist~\cite{eysenbach2004improving} is a
usable reporting standard, and Kitchenham and Pfleeger's series remains the
software engineering reference~\cite{kitchenham2008personal}.

\begin{pitfall}[The survey that measured nothing]
The characteristic failure: 300 responses, a 4.1 out of 5 mean agreement that
``documentation is important'', and no finding.

\medskip
It happens because the items were written without a research question behind
them, so no answer could have been surprising. Before writing any item, apply
the null test from \chref{experiments}: what distribution of answers would have
made you conclude the opposite? If no distribution would have, delete the item.
\end{pitfall}

\begin{traditions}{the survey}
\tradEMP{Estimates a population quantity with quantified uncertainty --- when
probability sampling is possible. When it is not, the survey is descriptive and
the write-up must say so.}
\tradDSR{Surveys practitioners to characterise the problem before designing,
and to gather perceived-utility evidence afterwards --- which is formative
naturalistic evaluation, and weak as summative evidence.}
\tradQUAL{Open-ended responses are qualitative data and deserve real analysis
(\chref{qda}), not a word cloud. Often the most informative part of the
instrument.}
\tradFORM{Little role, except that survey sampling theory is itself a formal
subject and design effects are computed, not guessed.}
\end{traditions}

\begin{lab}[Build and cognitively test an instrument]
\begin{tightnum}
  \item List your constructs. For each, decide: validated instrument, or your
    own items?
  \item Write items. Multi-item scales for anything continuous.
  \item Work template C8 end to end.
  \item Run cognitive interviews with three people from the target population.
    Record every item that was misread.
  \item Revise, then pilot with 10--15 people and run your full analysis script
    on the pilot data.
  \item Write the population, frame and coverage paragraph --- including who is
    excluded and in which direction it biases each main result.
\end{tightnum}

\medskip
This is portfolio piece P4.
\end{lab}

\begin{checkpoint}
\begin{tightnum}
  \item Distinguish population, sampling frame and achieved sample for a survey
    of Pakistani software professionals. Where is the largest error likely to
    be?
  \item Why can a 70\% response rate be more biased than a 20\% one?
  \item Rewrite: ``Do you agree that our new CI pipeline has made builds faster
    and more reliable?''
  \item You need to treat ``perceived usefulness'' as continuous in a
    regression. What does that require of your instrument design?
\end{tightnum}
\end{checkpoint}

\begin{chaptersummary}
\begin{tight}
  \item Surveys measure reported attitudes and practices at one time. They do
    not establish causation and do not measure behaviour.
  \item Population, frame and achieved sample are three things. Coverage error
    usually dominates sampling error in our field, and the confidence interval
    does not know about it.
  \item Non-probability sampling is normal and acceptable; claiming
    representativeness for it is not.
  \item Nonresponse bias, not response rate, is what matters. Check it with
    wave analysis, frame comparison or follow-up.
  \item Six wording faults account for most bad items. Behaviourally anchor
    anything socially loaded.
  \item Several items per construct if you want to treat it as continuous.
    Adopt validated instruments unmodified where they exist.
  \item Cognitive interviews are the highest-value hour in survey design.
  \item Report the funnel, the profile and the instrument in full.
\end{tight}
\end{chaptersummary}

\begin{reviewq}
  \item No sampling frame exists for ``software developers''. Given that, is
    probability sampling ever achievable in software engineering survey
    research? If so, describe a setting where it is; if not, say what the field
    should do about its collective reliance on convenience samples.
  \item Self-report and trace data on the same behaviour frequently disagree.
    Design a study that would establish the size and direction of that gap for
    one specific practice, and say what it would be worth knowing.
  \item Adopting a validated instrument imports its construct definition, which
    may have been developed in a different culture and language. Discuss the
    trade-off for research in Pakistan, and describe what re-validation would
    require.
  \item \reg{PhD Regs 2024}{\S14.3(B)} asks for external validity in both
    ecological and population senses. For a convenience-sampled survey, what is
    the most defensible thing a dissertation can claim, and how should the
    threats section be written?
\end{reviewq}
